\documentclass[11pt]{article}

\usepackage[margin=1in]{geometry}
\usepackage[T1]{fontenc}
\usepackage[utf8]{inputenc}
\usepackage{lmodern}
\usepackage{microtype}
\usepackage{amsmath,amssymb,amsthm,mathtools,bm}
\usepackage{enumitem}
\usepackage{booktabs}
\usepackage{longtable}
\usepackage{array}
\usepackage{xcolor}
\usepackage{hyperref}
\usepackage[nameinlink,noabbrev]{cleveref}

\hypersetup{
  colorlinks=true,
  linkcolor=blue!60!black,
  citecolor=blue!60!black,
  urlcolor=blue!60!black,
  pdftitle={A Control-Anchored Mean-Field Neural Differential Equation Formulation for P4-to-P60 Perturb-seq Dynamics},
  pdfauthor={OpenAI}
}

\title{A Control-Anchored Mean-Field Neural Differential Equation Formulation\\for P4-to-P60 Perturb-seq Dynamics}
\author{}
\date{}

\newcommand{\R}{\mathbb{R}}
\newcommand{\E}{\mathbb{E}}
\newcommand{\Prob}{\mathbb{P}}
\newcommand{\dd}{\,\mathrm{d}}
\newcommand{\KL}{\mathrm{KL}}
\newcommand{\UOT}{\mathrm{UOT}}
\newcommand{\SDunb}{\mathrm{SD}^{\mathrm{unb}}}
\newcommand{\diag}{\mathrm{diag}}
\newcommand{\tr}{\mathrm{tr}}
\newcommand{\softplus}{\mathrm{softplus}}

\newcommand{\Gcal}{\mathcal{G}}
\newcommand{\Gctrl}{\mathcal{G}_{\mathrm{ctrl}}}
\newcommand{\Ccal}{\mathcal{C}}
\newcommand{\Zcal}{\mathcal{Z}}
\newcommand{\simplex}{\Delta}
\newcommand{\Mplus}{\mathcal{M}_{+}}

\newcommand{\norm}[1]{\left\lVert #1 \right\rVert}
\newcommand{\abs}[1]{\left\lvert #1 \right\rvert}
\newcommand{\inner}[2]{\left\langle #1, #2 \right\rangle}

\begin{document}
\maketitle

\begin{abstract}
We formulate the $P4 \to P60$ perturb-seq problem as inference of a perturbation-indexed family of finite-measure-valued stochastic processes on a latent cell-state space. The biological motivation is that a perturbation can alter at least three coupled aspects of tumour evolution: directional progression through cell-state space, stochastic heterogeneity around that progression, and net expansion or depletion of the perturbed population. Because single-cell RNA sequencing is destructive, the observed data are snapshot distributions rather than tracked single-cell trajectories. This motivates a weighted stochastic differential equation formulation in a shared latent space. For each perturbation $g \in \Gcal$, the observed $P4$ cells and guide-abundance scale define an initial finite measure $\mu_0^g=M_g^0p_0^g$, while the observed $P60$ cells and guide-abundance scale define a terminal finite measure $\nu_1^g=M_g^1p_1^g$. The model is control anchored: control perturbations are represented by the zero embedding and therefore follow the baseline continuation law of the same coefficient networks. The model is also mean-field coupled: all perturbation-specific populations jointly determine a low-dimensional tissue context, and this context feeds back into the perturbation-conditioned dynamics. Mathematically, the resulting system is a reaction--drift--diffusion McKean--Vlasov model in latent space, implemented as a weighted neural stochastic differential equation with three separate control-anchored neural networks for directional movement, stochastic dispersion, and multiplicative growth or depletion. The predicted terminal finite measures are fitted to data with a single entropic unbalanced optimal-transport discrepancy, because terminal geometry and terminal abundance are two aspects of the same endpoint measure rather than two independent targets. Coarse cell-population labels may be used to build interpretable tissue-context summaries, but they are not required for the core endpoint supervision. The learned object is an effective interval generator for the single observed interval $P4 \to P60$; it is not claimed to be the unique microscopic biological law between the two snapshots.
\end{abstract}

\tableofcontents

\section{Biological background and modeling target}

\subsection{Biological setting}

The biological problem is to understand how a perturbation reshapes tumour evolution over the observed interval
\[
P4 \longrightarrow P60.
\]
At a biological level, a perturbation can affect at least three coupled mechanisms.

First, it can change the \emph{direction} in which cells move through transcriptomic state space. A perturbation may promote progression toward one epithelial program, block another program, or redirect cells into a different branch of the state manifold.

Second, it can change the amount of \emph{heterogeneity} or unresolved stochastic variation around that dominant direction. Even if two perturbations have similar average trends, one may produce a broader spread of cell states and another may confine the population more tightly.

Third, it can change the \emph{net abundance} of the perturbed population. A perturbation may confer a selective advantage, cause depletion, alter survival, or change proliferative capacity.

These three effects---movement, stochastic spread, and net growth or loss---are precisely the three components that motivate the mathematical structure of the model below.

\subsection{Why the data imply a population model rather than a trajectory-matching model}

Single-cell RNA sequencing is destructive. One does not observe the same physical cell at $P4$ and again at $P60$. Therefore, the data do not provide directly tracked cell lineages in state space. Instead, for each perturbation, one observes:
\begin{itemize}[leftmargin=2em]
    \item a collection of cell states at $P4$,
    \item a collection of cell states at $P60$,
    \item and guide-abundance information that provides a relative mass scale.
\end{itemize}

This means that the natural object of inference is not a set of matched cell-to-cell trajectories. The natural object is a \emph{population dynamics model} that evolves a perturbation-specific distribution or finite measure from the $P4$ snapshot to the $P60$ snapshot.

\subsection{Why control anchoring is natural}

The experimental design contains control perturbations, such as non-targeting guides. These controls provide a natural reference continuation law. Rather than learning a completely unrelated high-capacity dynamics for each perturbation, it is more statistically efficient and biologically interpretable to model each perturbation as a deviation from a shared control law. This is the reason for introducing a perturbation embedding $a_g$ with the exact anchor
\[
a_g=0
\qquad \text{for controls.}
\]

\subsection{Why mean-field tissue context is biologically plausible}

A perturbation can act cell-intrinsically, but the evolving tissue composition can also feed back on cell behavior. In a tumour setting, the current mixture of epithelial states, immune states, and signaling-related features can alter the effective continuation law. This motivates a mean-field interaction layer: the coefficient fields are allowed to depend not only on the current cell state and perturbation identity, but also on low-dimensional summaries of the current global population.

\subsection{What the model should estimate}

The target is a perturbation-conditioned family of evolving finite measures
\[
\Big\{ \mu_t^g : t\in[0,1],\ g\in\Gcal \Big\},
\]
where $t=0$ corresponds to $P4$ and $t=1$ corresponds to $P60$. The same model should explain both:
\begin{enumerate}[label=(\roman*),leftmargin=2.2em]
    \item where the perturbed population ends up in latent state space, and
    \item how much total perturbation-specific mass survives, expands, or depletes.
\end{enumerate}

\section{Explicit introduction to the SDE}

\subsection{What an SDE is}

A stochastic differential equation describes how a random state evolves over infinitesimal time. If $Z_t\in\R^d$ denotes a latent cell state, a generic SDE has the form
\begin{equation}
\dd Z_t = b(Z_t,t)\,\dd t + \sigma(Z_t,t)\,\dd W_t,
\label{eq:generic_sde}
\end{equation}
where:
\begin{itemize}[leftmargin=2em]
    \item $b(Z_t,t)$ is the \emph{drift}, the deterministic infinitesimal direction of motion;
    \item $\sigma(Z_t,t)$ is the \emph{diffusion map}, which scales stochastic fluctuations;
    \item $W_t$ is a Brownian motion, the canonical continuous-time noise process.
\end{itemize}

Over a small time interval $\Delta t$, \cref{eq:generic_sde} can be read heuristically as
\begin{equation}
Z_{t+\Delta t}-Z_t
\approx
b(Z_t,t)\,\Delta t
+
\sigma(Z_t,t)\sqrt{\Delta t}\,\xi_t,
\qquad
\xi_t \sim \mathcal{N}(0,I).
\label{eq:sde_increment}
\end{equation}
Thus, the drift term determines the average direction of motion, while the diffusion term determines the size and orientation of stochastic spread around that direction.

\subsection{Why an ODE alone is not enough}

A deterministic ODE would only include the drift term,
\[
\dd Z_t = b(Z_t,t)\,\dd t.
\]
That is too restrictive here. In the present problem, the latent state does not capture all biological variables, and the observed data consist only of sparse destructive snapshots. Therefore, a stochastic dispersion term is not merely a nuisance; it is part of the effective interval-level description. It absorbs unresolved heterogeneity, latent-state incompleteness, and intrinsically variable continuation around the main direction of motion.

\subsection{Why we also need a growth equation}

The SDE in \cref{eq:generic_sde} moves particles through latent state space, but it does not by itself account for perturbation-specific expansion or depletion. To model abundance changes, we attach a positive weight $w_t$ to each particle and write
\begin{equation}
\frac{\dd}{\dd t}\log w_t = r(Z_t,t),
\label{eq:generic_growth}
\end{equation}
where $r$ is a scalar growth or reaction field.

This has the explicit solution
\begin{equation}
w_t
=
\exp\!\left(
\int_0^t r(Z_s,s)\,\dd s
\right).
\label{eq:generic_growth_solution}
\end{equation}
If $r>0$ along the trajectory, the particle weight grows; if $r<0$, the weight shrinks.

\subsection{The weighted SDE we actually want}

For perturbation-conditioned tumour evolution, the natural conceptual model is therefore
\begin{equation}
\dd Z_t = b_g(Z_t,t;\mu_t)\,\dd t + \sigma_g(Z_t,t;\mu_t)\,\dd W_t,
\label{eq:conceptual_sde}
\end{equation}
together with
\begin{equation}
\frac{\dd}{\dd t}\log w_t = r_g(Z_t,t;\mu_t).
\label{eq:conceptual_weight}
\end{equation}

The perturbation index $g$ appears because different perturbations modify the continuation law. The dependence on $\mu_t$ appears because the current tissue context can feed back on the dynamics through a mean-field summary of the evolving population.

\subsection{From a single weighted particle to a population measure}

A single weighted particle does not represent the full perturbation-specific population. To represent the population, we initialize many particles at the observed $P4$ cells and then evolve them jointly. The perturbation-specific population at time $t$ is the weighted empirical measure
\begin{equation}
\mu_t^g
=
\frac{M_g^0}{n_g^0}
\sum_{i=1}^{n_g^0}
w_t^{g,i}\,\delta_{Z_t^{g,i}}.
\label{eq:conceptual_population_measure}
\end{equation}

This measure-valued object is the main mathematical prediction of the model. It contains both the terminal state allocation and the terminal perturbation mass.

\section{Notation table}

\renewcommand{\arraystretch}{1.16}

\begin{longtable}{p{0.16\textwidth}p{0.17\textwidth}p{0.59\textwidth}}
\caption{Core notation used throughout the formulation.}
\label{tab:notation}\\
\toprule
\textbf{Symbol} & \textbf{Type} & \textbf{Meaning} \\
\midrule
\endfirsthead

\toprule
\textbf{Symbol} & \textbf{Type} & \textbf{Meaning} \\
\midrule
\endhead

\bottomrule
\endfoot

$\Gcal$ & finite set & Set of perturbations.\\
$\Gctrl$ & subset of $\Gcal$ & Set of control perturbations.\\
$g,h$ & elements of $\Gcal$ & Perturbation indices.\\
$t\in[0,1]$ & scalar & Normalized time, with $t=0$ for $P4$ and $t=1$ for $P60$.\\
$\tau(t)$ & scalar & Optional map from normalized time to physical days.\\
$\Zcal\subseteq\R^d$ & space & Shared latent cell-state space.\\
$x_{gi},y_{gj}$ & vectors in $\Zcal$ & Observed latent states of cell $i$ at $P4$ and cell $j$ at $P60$ under perturbation $g$.\\
$X_g,Y_g$ & finite sets & Observed latent cell-state collections at $P4$ and $P60$ under perturbation $g$.\\
$n_g^0,n_g^1$ & integers & Numbers of observed $P4$ and $P60$ cells under perturbation $g$.\\
$p_0^g,p_1^g$ & probability measures & Empirical probability measures of the observed latent states at $P4$ and $P60$.\\
$M_g^0,M_g^1$ & positive scalars & Relative perturbation masses on the guide-abundance scale at $P4$ and $P60$.\\
$\mu_0^g,\nu_1^g$ & finite measures & Initial and observed terminal finite measures, $\mu_0^g=M_g^0p_0^g$ and $\nu_1^g=M_g^1p_1^g$.\\
$\mu_t^g$ & finite measure & Predicted perturbation-specific population measure at time $t$.\\
$\bar\mu_t$ & finite measure & Total population measure $\sum_{g\in\Gcal}\mu_t^g$.\\
$\bar M(t)$ & positive scalar & Total mass $\bar\mu_t(\Zcal)$.\\
$a_g$ & vector in $\R^r$ & Perturbation embedding; controls satisfy $a_g=0$.\\
$\omega$ & map $\Zcal\to\simplex^{K-1}$ & Optional coarse membership map used for interpretable context summaries.\\
$\phi$ & map $\Zcal\to\R^J$ & Optional mediator or signaling feature map.\\
$q_t$ & vector in $\simplex^{K-1}$ & Optional global coarse composition summary.\\
$s_t$ & vector in $\R^J$ & Optional global mediator summary.\\
$c_t$ & vector in $\R^L$ & Low-dimensional mean-field tissue-context vector.\\
$\gamma(t)$ & vector in $\R^{d_t}$ & Time embedding.\\
$u_t(z;\mu_t)$ & vector & Shared network input built from latent state, time, and tissue context.\\
$b_g$ & vector field & Drift field governing directional movement.\\
$\sigma_g$ & matrix field & Diffusion map governing stochastic dispersion.\\
$r_g$ & scalar field & Growth or reaction field governing multiplicative expansion or depletion.\\
$W_t^{g,i}$ & Brownian motion & Stochastic forcing for particle $i$ under perturbation $g$.\\
$Z_t^{g,i}$ & stochastic process & Latent state trajectory of particle $i$ under perturbation $g$.\\
$w_t^{g,i}$ & positive process & Weight of particle $i$ under perturbation $g$.\\
$\hat M_g^1$ & scalar & Predicted terminal mass $\mu_1^g(\Zcal)$.\\
$\hat p_1^g$ & probability measure & Predicted terminal normalized state distribution $\mu_1^g/\hat M_g^1$.\\
$\hat q_{g,1}$ & vector in $\simplex^{K-1}$ & Optional predicted terminal coarse composition.\\
$\sigma_{\min}$ & positive scalar & Lower bound used in the diffusion parameterization.\\
$r_{\max}$ & positive scalar & Bound used in the growth parameterization.\\
$\varepsilon$ & positive scalar & Entropic regularization parameter in unbalanced OT.\\
$\tau$ & positive scalar & Marginal-relaxation parameter in unbalanced OT.\\
$\UOT_{\varepsilon,\tau}$ & functional & Entropically regularized unbalanced OT cost on finite measures.\\
$\SDunb_{\varepsilon,\tau}$ & functional & Debiased unbalanced Sinkhorn divergence on finite measures.\\
$\theta_b,\theta_\sigma,\theta_r,\theta_c$ & parameters & Parameters of the drift, dispersion, growth, and context networks.\\
\end{longtable}

\section{Observed data and endpoint object}

\subsection{Normalized time}

We identify
\[
t=0 \equiv P4,
\qquad
t=1 \equiv P60.
\]
If one wishes to connect normalized time to physical days, one may set
\[
\tau(t)=4+56t,
\]
so that $\tau(0)=4$ and $\tau(1)=60$.

\subsection{Latent state space}

Let
\[
\Zcal \subseteq \R^d
\]
denote a common latent cell-state space obtained from pooled $P4$ and $P60$ transcriptomes. This latent space may be PCA, scVI, or another shared representation. The working assumption is that $\Zcal$ is sufficiently informative for an effective interval-level dynamics on $[0,1]$.

\subsection{Perturbation-indexed single-cell observations}

For each perturbation $g\in\Gcal$, let
\[
X_g = \{x_{g1},\dots,x_{g n_g^0}\}\subset \Zcal
\]
be the latent states of observed $P4$ cells carrying perturbation $g$, and let
\[
Y_g = \{y_{g1},\dots,y_{g n_g^1}\}\subset \Zcal
\]
be the latent states of observed $P60$ cells carrying perturbation $g$.

Define the empirical probability measures
\begin{equation}
p_0^g
=
\frac{1}{n_g^0}\sum_{i=1}^{n_g^0}\delta_{x_{gi}},
\qquad
p_1^g
=
\frac{1}{n_g^1}\sum_{j=1}^{n_g^1}\delta_{y_{gj}}.
\label{eq:empirical_prob_measures}
\end{equation}

\subsection{Guide-abundance masses and finite endpoint measures}

The numbers of sequenced cells and the guide-abundance measurements need not be on the same scale. Accordingly, we distinguish:
\begin{itemize}[leftmargin=2em]
    \item \emph{state samples}: $n_g^0$ and $n_g^1$, used to define $p_0^g$ and $p_1^g$;
    \item \emph{guide-abundance masses}: $M_g^0>0$ and $M_g^1>0$, representing relative perturbation abundance at $P4$ and $P60$ on the guide-abundance scale.
\end{itemize}

The initial and observed terminal finite measures are
\begin{equation}
\mu_0^g = M_g^0 p_0^g,
\qquad
\nu_1^g = M_g^1 p_1^g.
\label{eq:mass_scaled_measures}
\end{equation}

This is the key endpoint object. For each perturbation $g$, the observed terminal target is the \emph{finite measure} $\nu_1^g$, which already combines:
\begin{enumerate}[label=(\roman*),leftmargin=2.2em]
    \item terminal geometry in latent state space through $p_1^g$, and
    \item terminal abundance on the guide-abundance scale through $M_g^1$.
\end{enumerate}

\subsection{Optional interpretable tissue summaries}

If one wants a biologically interpretable mean-field context, one may define:
\[
\omega : \Zcal \to \simplex^{K-1},
\qquad
\phi : \Zcal \to \R^J,
\]
where $\omega$ is an optional coarse population membership map and $\phi$ is an optional mediator or signaling feature map.

For each perturbation $g$, the observed terminal coarse composition is then
\begin{equation}
q_{g,1}^{\mathrm{obs}}
=
\frac{1}{n_g^1}\sum_{j=1}^{n_g^1}\omega(y_{gj})
=
\frac{1}{M_g^1}\int_{\Zcal}\omega(z)\,\nu_1^g(\dd z).
\label{eq:observed_composition}
\end{equation}

These summaries are useful for interpretation, but they are not required for the core endpoint loss. The core supervision acts directly on finite measures.

\section{Assumptions}

The formulation rests on the following explicit assumptions.

\begin{enumerate}[label=(A\arabic*),leftmargin=2.8em]
    \item \textbf{Latent Markov sufficiency.}
    The pooled latent representation $\Zcal$ is sufficiently informative that the $P4\to P60$ interval can be described by an effective Markovian dynamics on $\Zcal$.

    \item \textbf{Within-perturbation exchangeability.}
    For each perturbation $g$, conditional on latent state and the shared tissue context, the joint law of cells carrying $g$ is invariant under permutation of cell indices. Thus, cells with the same perturbation are modeled as interchangeable particles, and the perturbation-specific empirical measure is the sufficient population-level state.

    \item \textbf{Relative-mass interpretation.}
    The masses $M_g^0$ and $M_g^1$ are treated as guide-abundance-scale masses and are therefore comparable across perturbations after normalization, even though they are distinct from the number of sequenced single-cell transcriptomes.

    \item \textbf{Mean-field closure.}
    The effect of the global tissue context on any cell is approximated through low-dimensional aggregate summaries of the joint population rather than through explicit pairwise interactions between all individual cells.

    \item \textbf{Control anchoring.}
    Control perturbations define the baseline continuation law. Perturbation effects are represented as low-dimensional deviations away from that baseline rather than as fully unrelated dynamics for each perturbation.

    \item \textbf{Effective-interval identifiability.}
    Because only two state-resolved time points are used, the learned dynamics are interpreted as an effective generator on $[0,1]$, not as a uniquely identified microscopic law at all intermediate times.
\end{enumerate}

\section{Control anchoring and mean-field tissue context}

\subsection{Control perturbations and perturbation embeddings}

Let
\[
\Gctrl \subset \Gcal
\]
denote the set of control perturbations. Each perturbation $g\in\Gcal$ is assigned an embedding
\[
a_g \in \R^r.
\]
To make the control anchor exact, we impose
\begin{equation}
a_g = 0
\qquad
\text{for every } g\in\Gctrl.
\label{eq:control_zero}
\end{equation}

If several sgRNAs target the same gene, one may optionally introduce
\begin{equation}
a_{gr} = a_g + \delta_{gr},
\qquad
\delta_{gr}\in\R^r,
\label{eq:guide_hierarchy}
\end{equation}
with shrinkage of $\delta_{gr}$ toward zero. For the core formulation below, a single perturbation-level embedding $a_g$ is sufficient.

\subsection{Joint population and total mass}

Let
\[
\mu_t = (\mu_t^g)_{g\in\Gcal}
\]
denote the family of perturbation-indexed populations at time $t$. The total population measure is
\begin{equation}
\bar\mu_t = \sum_{g\in\Gcal}\mu_t^g,
\qquad
\bar M(t)=\bar\mu_t(\Zcal).
\label{eq:global_population}
\end{equation}

\subsection{Generic mean-field context}

The most general formulation uses a low-dimensional tissue-context vector
\begin{equation}
c_t = \Psi_{\theta_c}(\mu_t)\in\R^L,
\label{eq:generic_context}
\end{equation}
where $\Psi_{\theta_c}$ is a context map from the current joint population law to a low-dimensional summary.

This expression is the precise mathematical content of the mean-field coupling: the coefficient fields are allowed to depend on the current population only through the context vector $c_t$.

\subsection{One interpretable choice of context map}

One interpretable implementation is to define
\begin{equation}
q_t
=
\frac{1}{\bar M(t)}
\sum_{g\in\Gcal}
\int_{\Zcal}\omega(z)\,\mu_t^g(\dd z)
\in \simplex^{K-1},
\label{eq:q_t}
\end{equation}
and
\begin{equation}
s_t
=
\frac{1}{\bar M(t)}
\sum_{g\in\Gcal}
\int_{\Zcal}\phi(z)\,\mu_t^g(\dd z)
\in \R^J.
\label{eq:s_t}
\end{equation}
Then set
\begin{equation}
c_t = \Psi_{\theta_c}(q_t,s_t)\in\R^L.
\label{eq:interpretable_context}
\end{equation}

This implementation keeps the mean-field layer interpretable in terms of coarse population composition and mediator features, but the core formulation does not require this particular choice.

\subsection{Time embedding and common network input}

Let
\[
\gamma : [0,1] \to \R^{d_t}
\]
be a time embedding. For a latent state $z\in\Zcal$, define
\begin{equation}
u_t(z;\mu_t)
=
\big(z,\ \gamma(t),\ c_t\big).
\label{eq:context_vector}
\end{equation}

This is the common input to the coefficient networks.

\section{Explicit control-anchored three-network weighted SDE}

\subsection{Why three separate networks}

The conceptual model in \cref{eq:conceptual_sde,eq:conceptual_weight} has three biologically distinct mechanisms:
\begin{enumerate}[label=(\roman*),leftmargin=2.2em]
    \item directional movement,
    \item stochastic dispersion,
    \item multiplicative growth or depletion.
\end{enumerate}
These mechanisms have different mathematical constraints and different identifiability properties. For that reason, they should not be collapsed into a single undifferentiated neural output. We therefore use:
\begin{enumerate}[label=(\arabic*),leftmargin=2.2em]
    \item a \emph{drift network},
    \item a \emph{dispersion network},
    \item a \emph{growth network}.
\end{enumerate}

\subsection{Drift network: directional movement}

Let
\[
\mathrm{DriftNet}_{\theta_b} :
\R^{d+d_t+L}
\to
\R^d \times \R^{d\times r},
\]
and write
\[
\mathrm{DriftNet}_{\theta_b}(u)=\big(\beta_b(u),B_b(u)\big),
\]
where
\[
\beta_b(u)\in\R^d,
\qquad
B_b(u)\in\R^{d\times r}.
\]
The perturbation-conditioned drift field is
\begin{equation}
b_g(z,t;\mu_t)
=
\beta_b\!\big(u_t(z;\mu_t)\big)
+
B_b\!\big(u_t(z;\mu_t)\big)\,a_g.
\label{eq:drift_total}
\end{equation}

The first term is the baseline control law; the second term is the perturbation-specific deviation from control.

\subsection{Dispersion network: stochastic spread}

Let
\[
\mathrm{DispersionNet}_{\theta_\sigma} :
\R^{d+d_t+L}
\to
\R^d \times \R^{d\times r},
\]
and write
\[
\mathrm{DispersionNet}_{\theta_\sigma}(u)=\big(\beta_\sigma(u),B_\sigma(u)\big),
\]
where
\[
\beta_\sigma(u)\in\R^d,
\qquad
B_\sigma(u)\in\R^{d\times r}.
\]

Define the pre-activation
\begin{equation}
\ell_g(z,t;\mu_t)
=
\beta_\sigma\!\big(u_t(z;\mu_t)\big)
+
B_\sigma\!\big(u_t(z;\mu_t)\big)\,a_g
\in\R^d,
\label{eq:disp_preact}
\end{equation}
and then define the diffusion matrix by
\begin{equation}
\sigma_g(z,t;\mu_t)
=
\diag\!\Big(
\softplus\big(\ell_g(z,t;\mu_t)\big)+\sigma_{\min}
\Big),
\qquad
\sigma_{\min}>0.
\label{eq:diff_total}
\end{equation}

The diagonal form is the default because diffusion is only weakly identifiable from two endpoint snapshots. If additional information becomes available, one may replace this with a low-rank or full-covariance parameterization.

\subsection{Growth network: expansion or depletion}

Let
\[
\mathrm{GrowthNet}_{\theta_r} :
\R^{d+d_t+L}
\to
\R \times \R^{1\times r},
\]
and write
\[
\mathrm{GrowthNet}_{\theta_r}(u)=\big(\beta_r(u),B_r(u)\big),
\]
where
\[
\beta_r(u)\in\R,
\qquad
B_r(u)\in\R^{1\times r}.
\]

Define the unconstrained log-rate
\begin{equation}
\tilde r_g(z,t;\mu_t)
=
\beta_r\!\big(u_t(z;\mu_t)\big)
+
B_r\!\big(u_t(z;\mu_t)\big)\,a_g,
\label{eq:growth_preact}
\end{equation}
and then set
\begin{equation}
r_g(z,t;\mu_t)
=
r_{\max}\tanh\!\big(\tilde r_g(z,t;\mu_t)\big),
\qquad
r_{\max}>0.
\label{eq:growth_total}
\end{equation}

The bounded $\tanh$ form stabilizes training because the particle weights depend exponentially on the time integral of $r_g$.

\subsection{Exact control anchor}

Because controls satisfy $a_g=0$, \cref{eq:drift_total,eq:disp_preact,eq:growth_preact} imply that control perturbations follow the baseline continuation law produced by
\[
\beta_b(\cdot),\qquad
\beta_\sigma(\cdot),\qquad
\beta_r(\cdot).
\]

\subsection{Particle system}

For each perturbation $g$ and each observed $P4$ cell $x_{gi}$, define a stochastic trajectory $Z_t^{g,i}$ and a positive weight $w_t^{g,i}$ by
\begin{equation}
\dd Z_t^{g,i}
=
b_g\big(Z_t^{g,i},t;\mu_t\big)\,\dd t
+
\sigma_g\big(Z_t^{g,i},t;\mu_t\big)\,\dd W_t^{g,i},
\qquad
Z_0^{g,i}=x_{gi},
\label{eq:particle_sde}
\end{equation}
\begin{equation}
\frac{\dd}{\dd t}\log w_t^{g,i}
=
r_g\big(Z_t^{g,i},t;\mu_t\big),
\qquad
w_0^{g,i}=1.
\label{eq:weight_ode}
\end{equation}
Here $W_t^{g,i}$ is a $d$-dimensional Brownian motion.

Therefore
\begin{equation}
w_t^{g,i}
=
\exp\!\left(
\int_0^t r_g\big(Z_s^{g,i},s;\mu_s\big)\,\dd s
\right).
\label{eq:weight_solution}
\end{equation}

\subsection{Population measure induced by the particle system}

The perturbation-specific population measure is
\begin{equation}
\mu_t^g
=
\frac{M_g^0}{n_g^0}
\sum_{i=1}^{n_g^0}
w_t^{g,i}\,\delta_{Z_t^{g,i}}
\in\Mplus(\Zcal).
\label{eq:mu_tg}
\end{equation}
In particular, the predicted terminal mass is
\begin{equation}
\hat M_g^1
=
\mu_1^g(\Zcal)
=
\frac{M_g^0}{n_g^0}
\sum_{i=1}^{n_g^0}w_1^{g,i}.
\label{eq:pred_terminal_mass}
\end{equation}

This weighted empirical measure is the central predicted object.

\subsection{Weak measure form}

For every smooth test function $f:\Zcal\to\R$,
\begin{equation}
\frac{\dd}{\dd t}\int_{\Zcal}f(z)\,\mu_t^g(\dd z)
=
\int_{\Zcal}
\left[
b_g(z,t;\mu_t)\cdot \nabla f(z)
+
\frac{1}{2}\tr\!\left(D_g(z,t;\mu_t)\nabla^2 f(z)\right)
+
r_g(z,t;\mu_t)f(z)
\right]
\mu_t^g(\dd z),
\label{eq:weak_form}
\end{equation}
where
\[
D_g(z,t;\mu_t)=\sigma_g(z,t;\mu_t)\sigma_g(z,t;\mu_t)^\top.
\]

If a density exists, \cref{eq:weak_form} corresponds formally to
\begin{equation}
\partial_t \mu_t^g
=
-\nabla\cdot\!\big(b_g(\cdot,t;\mu_t)\mu_t^g\big)
+
\frac{1}{2}\nabla\cdot\nabla:\!\big(D_g(\cdot,t;\mu_t)\mu_t^g\big)
+
r_g(\cdot,t;\mu_t)\mu_t^g.
\label{eq:fokker_planck}
\end{equation}

\subsection{Regularity conditions}

For well-posedness of the particle system, we assume:
\begin{enumerate}[label=(R\arabic*),leftmargin=2.8em]
    \item the outputs of the three networks are measurable in $(t,\mu)$ and locally Lipschitz in $z$;
    \item $\sigma_{\min}>0$ and $r_{\max}<\infty$ are fixed constants;
    \item whenever interpretable summaries $\omega$ and $\phi$ are used, they are bounded.
\end{enumerate}

\section{Explicit details behind the loss function}

\subsection{Why the loss must act on finite measures}

The forward model predicts, for each perturbation $g$, one terminal finite measure
\[
\mu_1^g.
\]
The observed endpoint is also one finite measure
\[
\nu_1^g = M_g^1 p_1^g.
\]
Therefore the statistically coherent supervision compares
\[
\mu_1^g
\qquad \text{to} \qquad
\nu_1^g
\]
directly as finite measures.

If one instead splits supervision into:
\begin{itemize}[leftmargin=2em]
    \item a normalized state-distribution loss on $\hat p_1^g$, and
    \item a separate scalar loss on $\hat M_g^1$,
\end{itemize}
one artificially decomposes one endpoint object into two losses. In the present formulation, terminal geometry and terminal abundance are fitted jointly by a single unbalanced discrepancy.

\subsection{Derived endpoint summaries}

Although the primary predicted object is $\mu_1^g$, it is often useful to report
\begin{equation}
\hat M_g^1 = \mu_1^g(\Zcal),
\qquad
\hat p_1^g = \frac{\mu_1^g}{\hat M_g^1},
\label{eq:pred_terminal_prob}
\end{equation}
and, when $\omega$ is used,
\begin{equation}
\hat q_{g,1}
=
\frac{1}{\hat M_g^1}
\int_{\Zcal}\omega(z)\,\mu_1^g(\dd z)
\in \simplex^{K-1}.
\label{eq:pred_terminal_comp}
\end{equation}
These are summaries of the finite measure, not replacements for it.

\subsection{Generalized KL divergence for finite measures}

For finite nonnegative measures $\alpha,\beta\in\Mplus(\Zcal)$, define the generalized KL divergence by
\begin{equation}
\KL(\alpha \,\|\, \beta)
=
\begin{cases}
\displaystyle
\int_{\Zcal}
\log\!\left(\frac{\dd\alpha}{\dd\beta}\right)\,\alpha(\dd z)
-
\alpha(\Zcal)
+
\beta(\Zcal),
& \alpha \ll \beta, \\[1.5ex]
+\infty, & \text{otherwise.}
\end{cases}
\label{eq:gkl}
\end{equation}
This is the correct KL notion when the compared objects need not have the same total mass.

\subsection{Entropic unbalanced optimal transport}

Let $c(z,z')$ be the latent ground cost, typically
\[
c(z,z')=\norm{z-z'}_2^2.
\]
For finite measures $\alpha,\beta\in\Mplus(\Zcal)$, define
\begin{equation}
\UOT_{\varepsilon,\tau}(\alpha,\beta)
=
\min_{\pi\in\Mplus(\Zcal\times\Zcal)}
\left\{
\int_{\Zcal\times\Zcal}c(z,z')\,\pi(\dd z,\dd z')
+
\varepsilon\,\KL(\pi\,\|\,\alpha\otimes\beta)
+
\tau\,\KL(\pi_1\,\|\,\alpha)
+
\tau\,\KL(\pi_2\,\|\,\beta)
\right\},
\label{eq:uot_def}
\end{equation}
where $\pi_1,\pi_2$ are the marginals of $\pi$.

This objective has four parts:
\begin{enumerate}[label=(\roman*),leftmargin=2.2em]
    \item the transport term penalizes moving mass too far in latent space;
    \item the entropic term stabilizes the problem numerically;
    \item the marginal KL terms allow creation or destruction of mass relative to strict balanced transport;
    \item the parameter $\tau$ controls how costly that mass variation is.
\end{enumerate}

\subsection{Debiased unbalanced Sinkhorn divergence}

To reduce entropic bias, define
\begin{equation}
\SDunb_{\varepsilon,\tau}(\alpha,\beta)
=
\UOT_{\varepsilon,\tau}(\alpha,\beta)
-
\frac{1}{2}\UOT_{\varepsilon,\tau}(\alpha,\alpha)
-
\frac{1}{2}\UOT_{\varepsilon,\tau}(\beta,\beta).
\label{eq:unbalanced_sinkhorn}
\end{equation}

\subsection{Primary endpoint loss}

The primary endpoint loss is
\begin{equation}
\mathcal{L}_{\mathrm{end}}
=
\sum_{g\in\Gcal}
\alpha_g^{\mathrm{end}}
\,\SDunb_{\varepsilon,\tau}\!\left(\mu_1^g,\nu_1^g\right),
\label{eq:L_end}
\end{equation}
where $\alpha_g^{\mathrm{end}}$ are user-chosen weights.

This loss is the central reason to formulate the endpoint as a finite-measure problem. It jointly penalizes:
\begin{itemize}[leftmargin=2em]
    \item mismatch in terminal state geometry,
    \item mismatch in total terminal mass,
    \item and mismatch in the tradeoff between latent movement and mass change.
\end{itemize}

\subsection{Optional auxiliary composition diagnostic}

If one wants a weak interpretable diagnostic, one may define
\begin{equation}
\mathcal{L}_{\mathrm{comp}}^{\mathrm{aux}}
=
\sum_{g\in\Gcal}
\alpha_g^{\mathrm{comp}}
\,\KL\!\left(q_{g,1}^{\mathrm{obs}}\,\|\,\hat q_{g,1}\right).
\label{eq:L_comp_aux}
\end{equation}
Because $q_{g,1}^{\mathrm{obs}}$ is derived from the same endpoint single-cell data that define $p_1^g$, this term is optional and not part of the primary supervision.

\section{Regularization and identifiability}

\subsection{Embedding shrinkage}

For non-control perturbations,
\begin{equation}
\mathcal{R}_{\mathrm{emb}}
=
\sum_{g\in\Gcal\setminus\Gctrl}\norm{a_g}^2.
\label{eq:R_emb}
\end{equation}
If \cref{eq:guide_hierarchy} is used, add
\begin{equation}
\sum_g \sum_{r\in\mathcal{R}_g}\norm{\delta_{gr}}^2.
\label{eq:R_sgrna}
\end{equation}

\subsection{Perturbation-modulation regularization}

Let
\[
\tilde p_t^g = \frac{\mu_t^g}{\mu_t^g(\Zcal)}
\]
denote the normalized predicted population of perturbation $g$ at time $t$. A natural penalty on perturbation-specific deviations from control is
\begin{equation}
\mathcal{R}_{\mathrm{mod}}
=
\sum_{g\in\Gcal\setminus\Gctrl}
\int_0^1
\E_{z\sim \tilde p_t^g}
\left[
\norm{B_b(u_t(z;\mu_t))a_g}^2
+
\norm{B_\sigma(u_t(z;\mu_t))a_g}^2
+
\abs{B_r(u_t(z;\mu_t))a_g}^2
\right]
\dd t.
\label{eq:R_mod}
\end{equation}

\subsection{Dispersion regularization}

Because diffusion is only weakly identifiable from two endpoint snapshots, the model benefits from discouraging excessive stochastic spread:
\begin{equation}
\mathcal{R}_{\mathrm{disp}}
=
\sum_{g\in\Gcal}
\int_0^1
\E_{z\sim\tilde p_t^g}
\left[
\norm{\sigma_g(z,t;\mu_t)}_F^2
\right]
\dd t.
\label{eq:R_disp}
\end{equation}

\subsection{Network parameter regularization}

A generic parameter penalty is
\begin{equation}
\mathcal{R}_{\mathrm{nn}}
=
\norm{\theta_b}^2
+
\norm{\theta_\sigma}^2
+
\norm{\theta_r}^2
+
\norm{\theta_c}^2.
\label{eq:R_nn}
\end{equation}

\subsection{Why strong regularization is necessary}

From only two state-resolved time points, drift, diffusion, and growth are not equally identifiable. In practice:
\begin{itemize}[leftmargin=2em]
    \item drift carries the main burden of explaining large-scale state change;
    \item diffusion should remain relatively simple and strongly regularized;
    \item growth must be bounded because it exponentiates into the particle weights.
\end{itemize}
These are not merely numerical choices; they are part of the effective interval modeling assumptions.

\section{Full training objective}

The primary objective is
\begin{equation}
\mathcal{L}
=
\lambda_{\mathrm{end}}\mathcal{L}_{\mathrm{end}}
+
\lambda_{\mathrm{emb}}\mathcal{R}_{\mathrm{emb}}
+
\lambda_{\mathrm{mod}}\mathcal{R}_{\mathrm{mod}}
+
\lambda_{\mathrm{disp}}\mathcal{R}_{\mathrm{disp}}
+
\lambda_{\mathrm{nn}}\mathcal{R}_{\mathrm{nn}}
+
\lambda_{\mathrm{comp}}^{\mathrm{aux}}\mathcal{L}_{\mathrm{comp}}^{\mathrm{aux}}.
\label{eq:full_objective}
\end{equation}
The default main formulation takes
\[
\lambda_{\mathrm{comp}}^{\mathrm{aux}}=0
\]
or very small.

\section{Numerical approximation and training}

\subsection{Euler--Maruyama approximation}

Let
\[
t_\ell=\ell\Delta t,
\qquad
\Delta t=\frac{1}{L},
\qquad
\ell=0,\dots,L-1.
\]
Using \cref{eq:particle_sde,eq:weight_ode}, a standard Euler--Maruyama update is
\begin{equation}
Z_{\ell+1}^{g,i}
=
Z_{\ell}^{g,i}
+
b_g\!\left(Z_\ell^{g,i},t_\ell;\mu_{t_\ell}\right)\Delta t
+
\sigma_g\!\left(Z_\ell^{g,i},t_\ell;\mu_{t_\ell}\right)\sqrt{\Delta t}\,\xi_\ell^{g,i},
\label{eq:em_state}
\end{equation}
\begin{equation}
w_{\ell+1}^{g,i}
=
w_{\ell}^{g,i}
\exp\!\left(
r_g\!\left(Z_\ell^{g,i},t_\ell;\mu_{t_\ell}\right)\Delta t
\right),
\label{eq:em_weight}
\end{equation}
where
\[
\xi_\ell^{g,i}\sim\mathcal{N}(0,I_d).
\]

\subsection{Context update}

If the interpretable context construction \cref{eq:q_t,eq:s_t,eq:interpretable_context} is used, then at each step one recomputes
\begin{equation}
q_{t_\ell}
=
\frac{1}{\bar M(t_\ell)}
\sum_{g\in\Gcal}
\frac{M_g^0}{n_g^0}
\sum_{i=1}^{n_g^0}
w_\ell^{g,i}\,\omega(Z_\ell^{g,i}),
\label{eq:discrete_q}
\end{equation}
\begin{equation}
s_{t_\ell}
=
\frac{1}{\bar M(t_\ell)}
\sum_{g\in\Gcal}
\frac{M_g^0}{n_g^0}
\sum_{i=1}^{n_g^0}
w_\ell^{g,i}\,\phi(Z_\ell^{g,i}),
\label{eq:discrete_s}
\end{equation}
and then set
\begin{equation}
c_{t_\ell}=\Psi_{\theta_c}(q_{t_\ell},s_{t_\ell}).
\label{eq:discrete_c}
\end{equation}
If a different context map is used, one updates $c_{t_\ell}$ according to that construction instead.

\subsection{Discrete terminal measures}

At the terminal step $t_L=1$, the predicted finite measure is
\begin{equation}
\mu_1^g
=
\frac{M_g^0}{n_g^0}
\sum_{i=1}^{n_g^0}
w_L^{g,i}\,\delta_{Z_L^{g,i}},
\label{eq:mu_terminal_discrete}
\end{equation}
while the observed terminal measure is
\begin{equation}
\nu_1^g
=
\frac{M_g^1}{n_g^1}
\sum_{j=1}^{n_g^1}\delta_{y_{gj}}.
\label{eq:nu_terminal_discrete}
\end{equation}
The discrete unbalanced OT problem underlying \cref{eq:L_end} is then built from the cost matrix
\[
C_{ij}^g = c(Z_L^{g,i},y_{gj}).
\]

\subsection{Practical training loop}

A practical training loop is:
\begin{enumerate}[label=(\arabic*),leftmargin=2.2em]
    \item construct the shared latent space $\Zcal$ from pooled $P4$ and $P60$ data;
    \item initialize particles at the observed $P4$ cells for each perturbation;
    \item jointly simulate all perturbation-indexed populations from $t=0$ to $t=1$ using \cref{eq:em_state,eq:em_weight};
    \item compute the terminal measures $\mu_1^g$ and the derived summaries $\hat M_g^1$, $\hat p_1^g$, and optionally $\hat q_{g,1}$;
    \item evaluate the primary endpoint loss \cref{eq:L_end} and the regularizers \cref{eq:R_emb,eq:R_mod,eq:R_disp,eq:R_nn};
    \item optionally evaluate \cref{eq:L_comp_aux} for diagnosis or weak auxiliary regularization;
    \item update parameters by stochastic gradient descent.
\end{enumerate}

If jointly simulating all perturbations is too costly, one may use a stratified minibatch of perturbations to estimate the mean-field terms. That changes the numerical estimator, not the population-level model.

\section{Biological outputs and counterfactuals}

The formulation produces the following outputs.

\subsection{Predicted terminal finite measure}

The central endpoint object is
\[
\mu_1^g,
\]
which jointly encodes terminal abundance and terminal state allocation.

\subsection{Predicted abundance effect}

The abundance or selection effect of perturbation $g$ is summarized by
\[
\hat M_g^1
\quad \text{or} \quad
\log \hat M_g^1 - \log M_g^0.
\]

\subsection{Predicted terminal state distribution}

The normalized endpoint state distribution is
\[
\hat p_1^g.
\]

\subsection{Predicted terminal coarse composition}

If $\omega$ is used, the predicted coarse endpoint allocation is
\[
\hat q_{g,1}.
\]
This is an interpretive summary, not the main fitted target.

\subsection{Intrinsic versus context-mediated effect}

Because the model separates perturbation embedding from tissue context, one may compare:
\begin{enumerate}[label=(\alph*),leftmargin=2.2em]
    \item a full simulation using the self-consistent mean field $c_t$, and
    \item a simulation with $c_t$ clamped to a reference trajectory or removed.
\end{enumerate}
This gives a principled decomposition of endpoint effects into intrinsic perturbation components and context-mediated components.

\subsection{Control-referenced perturbation effect}

Because controls correspond to $a_g=0$, one may also compare:
\begin{enumerate}[label=(\alph*),leftmargin=2.2em]
    \item evolution under $a_g$, and
    \item evolution under $a=0$,
\end{enumerate}
starting from the same initial state distribution. This directly quantifies the perturbation effect relative to the control continuation law.

\section{What this model is, and is not}

\subsection{It is a mean-field interaction model}

Mathematically, the model is a perturbation-indexed McKean--Vlasov reaction--drift--diffusion system with low-dimensional population feedback through the tissue-context vector $c_t$.

\subsection{It is not a full mean-field game}

A literal mean-field game would require explicit control variables, value functions, and a coupled Hamilton--Jacobi--Bellman / Fokker--Planck system. None of those are introduced here. The model uses mean-field feedback, but it is not a full equilibrium game.

\subsection{It estimates an effective interval generator}

From two state-resolved time points, the continuous-time law between them is not uniquely identifiable. The learned coefficient fields must therefore be interpreted as a regularized effective generator on $[0,1]$, not as the unique microscopic biological truth.

\subsection{It supports intervention-conditional interpretation, not unrestricted causal claims}

The model is suitable for counterfactual simulation within the support of the observed perturbation-conditioned data and latent representation. It does not, by itself, identify unrestricted causal effects outside that support.

\section{Concise summary}

The formulation can be summarized compactly as follows:

\begin{quote}
For each perturbation $g$, start from the observed $P4$ finite measure $\mu_0^g=M_g^0p_0^g$, evolve a weighted particle system in latent space under a control-anchored perturbation-conditioned McKean--Vlasov neural SDE whose drift, stochastic dispersion, and growth are modeled by three separate neural networks conditioned on latent state, time, and low-dimensional tissue context, and fit the resulting terminal finite measure $\mu_1^g$ to the observed $P60$ finite measure $\nu_1^g=M_g^1p_1^g$ by a single unbalanced optimal-transport discrepancy.
\end{quote}

Equivalently, the model has five layers:
\begin{enumerate}[label=(\arabic*),leftmargin=2.2em]
    \item a \emph{biological mechanism layer}, which separates movement, stochastic spread, and growth/depletion;
    \item an \emph{SDE layer}, which evolves latent cell states;
    \item a \emph{weight layer}, which evolves perturbation-specific mass;
    \item a \emph{mean-field context layer}, which feeds back low-dimensional tissue information;
    \item an \emph{endpoint loss layer}, which compares predicted and observed terminal finite measures by unbalanced optimal transport.
\end{enumerate}

This is the explicit formal version of the $P4 \to P60$ control-anchored mean-field neural differential equation formulation after placing the biology first, explicitly introducing the SDE, and then giving the detailed finite-measure model and loss function.

\end{document}